\chapter*{Glosario de términos}
\addcontentsline{toc}{chapter}{Glosario de términos}
\label{cap:glosario}

Este glosario recoge, ordenados alfabéticamente, los términos técnicos y
las siglas usadas a lo largo de la memoria que pueden no resultar
familiares fuera del ámbito del seguimiento multiobjeto y la visión por
computador. Cada capítulo introduce y explica estos conceptos la primera
vez que aparece.

\begin{description}
\item[AssA (\emph{Association Accuracy})] Componente de HOTA que mide la
  calidad de la asociación entre detecciones de un mismo objeto a lo largo
  del tiempo, con independencia de si la detección en sí es correcta
  (sección~\ref{subsec:fundamentos-hota}).

\item[ByteTrack] Algoritmo de \emph{tracking-by-detection} que asocia
  también las detecciones de baja confianza en una segunda ronda, en vez
  de descartarlas directamente (sección~\ref{subsec:fundamentos-bytetrack}).

\item[\texttt{det.txt}] Fichero de detecciones precomputadas de una
  secuencia de SoccerNet-Tracking, en formato MOT: una fila por detección,
  con \emph{frame}, coordenadas de la caja delimitadora y confianza
  (sección~\ref{subsec:fundamentos-soccernet-tracking}).

\item[DetA (\emph{Detection Accuracy})] Componente de HOTA que mide la
  calidad de las detecciones por sí solas, sin tener en cuenta si están
  bien asociadas entre \emph{frames} (sección~\ref{subsec:fundamentos-hota}).

\item[\emph{Fourcc}] Código de cuatro caracteres que identifica el códec
  de vídeo usado al codificar un fichero con \texttt{OpenCV}; el proyecto
  usa \texttt{avc1} (H.264) en vez de \texttt{mp4v} para que los vídeos se
  reproduzcan en el navegador (problema P17,
  Tabla~\ref{tab:impl-problemas}).

\item[\emph{Ground truth}] Anotación de referencia considerada correcta,
  usada para evaluar un sistema; en SoccerNet-Tracking, el fichero
  \texttt{gt.txt} con las posiciones e identidades reales de cada jugador
  (sección~\ref{subsec:fundamentos-soccernet-tracking}).

\item[Homografía] Transformación geométrica que proyecta puntos de un
  plano (la imagen de la cámara) a otro plano (el campo real), usada en
  este proyecto para convertir posiciones en píxeles a metros reales
  (sección~\ref{sec:fundamentos-homografia}).

\item[HOTA (\emph{Higher Order Tracking Accuracy})] Métrica de evaluación
  de \emph{tracking} que combina DetA y AssA, promediada sobre un rango de
  umbrales de IoU (0.05 a 0.95); en este proyecto se distingue
  explícitamente de HOTA(0) (sección~\ref{subsec:fundamentos-hota},
  sección~\ref{subsec:impl-hota-vs-hota0}).

\item[HOTA(0)] Notación usada en esta memoria para el valor de HOTA
  evaluado en un único umbral de IoU (0.5), en vez del promedio integrado
  de la métrica HOTA estándar; la distinción entre ambas es uno de los
  hallazgos centrales del proyecto
  (sección~\ref{subsec:impl-hota-vs-hota0}).

\item[Húngaro, algoritmo] Algoritmo de optimización combinatoria que
  resuelve el problema de asignación óptima entre dos conjuntos, en este
  contexto, qué detección se asocia a qué \emph{track} en tiempo
  polinómico (sección~\ref{subsec:fundamentos-hungaro}).

\item[IDF1] Métrica de evaluación de \emph{tracking} centrada en la
  identidad: mide qué proporción de las detecciones se asocian al
  identificador correcto a lo largo de toda la secuencia
  (sección~\ref{sec:fundamentos-metricas}).

\item[IDSW (\emph{ID Switches})] Número de veces que un algoritmo de
  \emph{tracking} cambia por error el identificador asignado a un mismo
  objeto real (sección~\ref{sec:fundamentos-metricas}).

\item[IoU (\emph{Intersection over Union})] Medida de solapamiento entre
  dos cajas delimitadoras: el área de su intersección dividida entre el
  área de su unión, usada tanto para asociar detecciones como para
  evaluar \emph{trackers} (sección~\ref{subsec:fundamentos-iou}).

\item[K-means] Algoritmo de agrupamiento no supervisado que reparte un
  conjunto de puntos en $k$ grupos según su cercanía; usado en este
  proyecto para clasificar jugadores por equipo a partir del color de su
  camiseta (sección~\ref{sec:fundamentos-kmeans}).

\item[Kalman, filtro de] Algoritmo recursivo de estimación que predice el
  estado futuro de un sistema (aquí, la posición de un jugador) y lo
  corrige con cada nueva observación, base de los algoritmos de
  \emph{tracking-by-detection} usados en este proyecto
  (sección~\ref{subsec:fundamentos-kalman}).

\item[KDE (\emph{Kernel Density Estimation})] Técnica estadística para
  estimar la densidad de un conjunto de puntos de forma continua, usada
  para generar los mapas de calor de posiciones a partir de las
  trayectorias de \emph{tracking} (sección~\ref{sec:implementacion-metadatos}).

\item[MOT (\emph{Multi-Object Tracking})] Seguimiento multiobjeto: la
  tarea de detectar y seguir la identidad de múltiples objetos a la vez a
  lo largo de una secuencia de vídeo, tarea central de este proyecto
  (sección~\ref{sec:fundamentos-mot}).

\item[MOTA (\emph{Multiple Object Tracking Accuracy})] Métrica clásica de
  evaluación de \emph{tracking} que combina falsos positivos, falsos
  negativos y cambios de identidad en un único valor
  (sección~\ref{sec:fundamentos-metricas}).

\item[MVP (\emph{Minimum Viable Product})] Versión mínima funcional de un
  producto, usada aquí para describir el alcance de la aplicación web
  implementada: fiel al diseño planificado, sin todas las ampliaciones
  posibles (sección~\ref{sec:implementacion-web}).

\item[OC-SORT (\emph{Observation-Centric SORT})] Variante de SORT que
  reutiliza las detecciones originales sin interpolarlas y corrige la
  deriva del filtro de Kalman durante oclusiones prolongadas
  (sección~\ref{subsec:fundamentos-ocsort}).

\item[Re-ID (re-identificación)] Técnica que usa la apariencia visual de
  un objeto, además de su posición, para reasociar su identidad tras una
  oclusión; considerada y descartada para este proyecto
  (sección~\ref{subsec:fundamentos-reid}).

\item[\emph{Seqmap}] Fichero que indica a \texttt{TrackEval} qué
  secuencias evaluar en una ejecución
  (sección~\ref{sec:manual-evaluacion}).

\item[SORT (\emph{Simple Online and Realtime Tracking})] Algoritmo
  clásico de \emph{tracking-by-detection} que combina el filtro de Kalman
  con el algoritmo húngaro sobre IoU, base de la que parte OC-SORT
  (sección~\ref{subsec:fundamentos-sort}).

\item[SPA (\emph{Single Page Application})] Aplicación web que carga una
  única página HTML y actualiza su contenido dinámicamente en el
  navegador, sin recargas completas; arquitectura de la aplicación web de
  este proyecto (sección~\ref{sec:diseno-web}).

\item[SoccerNet-Tracking] \emph{Dataset} público de vídeo de fútbol con
  detecciones precomputadas y \emph{ground truth}, usado como base de
  datos de este proyecto
  (sección~\ref{subsec:fundamentos-soccernet-tracking}).

\item[\emph{Split}] Cada una de las particiones del \emph{dataset}
  (entrenamiento, test) sobre las que se ejecutan y evalúan los
  algoritmos por separado (sección~\ref{subsec:fundamentos-soccernet-tracking}).

\item[\emph{Tracking-by-detection}] Enfoque de seguimiento multiobjeto
  que parte de detecciones ya calculadas en cada \emph{frame} y se centra
  en asociarlas entre sí, en vez de detectar y seguir a la vez
  (sección~\ref{sec:fundamentos-tbd}).

\item[TrackEval] Librería de referencia usada en este proyecto para
  calcular HOTA, MOTA e IDF1 de forma estandarizada
  (sección~\ref{sec:implementacion-evaluacion}).

\item[\emph{Tracker}] Nombre genérico para un algoritmo de seguimiento
  multiobjeto; en este proyecto, ByteTrack u OC-SORT.
\end{description}
